% page 79 du fac-simile (image p-078 du scan)
% bloc 165.1 x 229.0 mm ; marge gauche 36.9 mm
\pgc{15.240mm}{0.000mm}{
\l{\cel{54}71.}
\l{}
\l{\sou{boto}. Pedo-vesto ledra qua kontenas la gambo. (\sou{Boteto}\cel{1}kontenas}
\l{\cel{3}nur la infra parto di la gambo ;\cel{1}\sou{shuo}\cel{1}kontenas nur la pedo).}
\l{\cel{3}- EFRS.}
\l{}
\l{\sou{botaniko}. Parto di la natur-historio qua havas kom temo di}
\l{\cel{3}studio la vejetanti. - DEFIRS.}
\l{}
\l{\sou{botanizar}. (\sou{netrans.}) Kolektar por la studio botanikala.}
\l{\cel{3}(ne \textquotedbl{}praktikar o studiar botaniko\textquotedbl{}). - DEFIRS.}
\l{}
\l{\sou{botelo}. Vazo portebla, di qua la ventro esas cilindra, la kolo}
\l{\cel{3}streta e longa, destinita a kontenar liquidi, precipue vino.}
\l{\cel{3}- DEFIRS.}
\l{}
\l{\sou{bovo}. (\sou{zool,}) Mamifero ruminera quan onu uzas prefere por la}
\l{\cel{3}agrokultivo se lu esas bovulo, e por obtenar lua lakto se lu}
\l{\cel{3}esas bovino - o quan onu grasigas por bucheso. (Kande bovulo}
\l{\cel{3}esas destinita a la tauro-kombati, lua nomo esas \textquotedbl{}tauro\textquotedbl{}). -}
\l{\cel{2}EFIS.}
\l{}
\l{\sou{boxar}.\cel{2}(\sou{trans. e netrans}.) Interbatar per pugno-frapi, en la}
\l{\cel{3}maniero Angla. - DEFR.}
\l{}
\l{\sou{boyo.}\cel{2}(\sou{navig}.)Omna korpo flotacanta (peco de ligno o korko,}
\l{\cel{3}barelo, e c.) destinita a markizar, sur la surfaco di aquo,}
\l{\cel{3}la punto ube lansesis ankro - ed informar pri rifo od obstaklo}
\l{\cel{3}od indikar la direciono di paseyo o di kanaleto. - DEFIRS.}
\l{}
\l{\sou{boyardo}. - Sinioro, personego, en Rusia caroza, en Transilvania.}
\l{\cel{3}- DEFIRS.}
\l{}
\l{\sou{braco}. Kordego an la du extremaji di yardo. - DEFIS.}
\l{}
\l{\sou{braceleto}. Ornivo ringo-forma qua cirkondas la ponyeto. - DEFIRS.}
\l{}
\l{\sou{bracho}. Viro-vesto, qua kovras la persono, de la hanchi til sub}
\l{\cel{3}la genui, e bifurkas tale ke olu envelopas la kruri. - EFIS.}
\l{}
\l{\sou{bradipo}. (\sou{zool.}) Mamifero de la tribuo \textquotedbl{}sen-denti\textquotedbl{}. L.\cel{1}\sou{bradypus}.}
\l{\cel{3}- EFIS.}
\l{}
\l{\sou{brako}. Speco di chas-hundo, reza-pila, e kun oreli pendanta.}
\l{\cel{3}- DEFIS.}
\l{}
\l{\sou{brakio}. I. Membro en la parto supra di la korpo homala, qua}
\l{\cel{3}artikizesas ye la shultro e qua finas per la manuo ; che ula}
\l{\cel{3}animali, membri qui korespondas a la brakii homala pro lia}
\l{\cel{3}situeso o funciono. - II. Parto de la maro qua konstriktesas}
\l{\cel{3}da du litori. - DEFIRS.}
\l{}
\l{\sou{brakicefala}. Di qua la kapo esas kurta. - DEFIRS.}
}
